\lysh{18865}{4}
\begin{alist}
\item Αφού το πεδίο ορισμού της συνάρτησης $f$ είναι το σύνολο όλων των πραγματικών αριθμών, για τους οποίους το $f(x)$ έχει νόημα, είναι:
\[|x| > 0 \Leftrightarrow x \ne 0\]
Επομένως $D_f = (-\infty, 0) \cup (0, +\infty)$.
\item Είναι: $f(-x) = \ln|-x| = \ln|x| = f(x)$. \\
Έτσι για κάθε $x \in D_f$ το $-x \in D_f$ και ισχύει $f(-x) = f(x)$. \\
Επομένως η συνάρτηση είναι άρτια και ως εκ τούτου συμμετρική ως προς τον άξονα $y'y$.
\item Είναι: 
\[f(x) = \ln|x| = \begin{cases} \ln x, & \text{αν } x > 0 \\ \ln(-x), & \text{αν } x < 0 \end{cases}\]
Η γραφική παράσταση της συνάρτησης $f$ αποτελείται από δύο κλάδους. \\
Αν $x > 0$, τότε έχουμε τη γραφική παράσταση της λογαριθμικής συνάρτησης $y = \ln x$. \\
Αν $x < 0$, παίρνουμε την συμμετρική καμπύλη της λογαριθμικής συνάρτησης $y = \ln x$ ως προς τον άξονα $y'y$. \\
Επομένως, η γραφική παράσταση της συνάρτησης $f$ είναι η ακόλουθη:
\begin{center}
\begin{tikzpicture}
\begin{axis}[
    axis lines=middle,
    xmin=-4.2, xmax=4.2,
    ymin=-3.5, ymax=1.5,
    xtick={-4, -3.5, -3, -2.5, -2, -1.5, -1, -0.5, 0.5, 1, 1.5, 2, 2.5, 3, 3.5, 4},
    ytick={-3, -2.5, -2, -1.5, -1, -0.5, 0.5, 1},
    xticklabels={-4, -3.5, -3, -2.5, -2, -1.5, -1, -0.5, 0.5, 1, 1.5, 2, 2.5, 3, 3.5, },
    yticklabels={-3, -2.5, -2, -1.5, -1, -0.5, 0.5, 1},
    grid=both,
    grid style={line width=.1pt, draw=gray!30},
    major grid style={line width=.2pt,draw=gray!50},
    thick,
    width=14cm, height=8cm,
    samples=200
]
\addplot[domain=0.03:4.2, maincolor, very thick] {ln(x)} node[pos=0.1, above left, text=black] {$\text{C}_f$};
\addplot[domain=-4.2:-0.03, maincolor, very thick] {ln(-x)};
\end{axis}
\end{tikzpicture}
\end{center}
\item Υπάρχουν δύο περιπτώσεις, όπως φαίνεται στο παρακάτω σχήμα:
\end{alist}

